\chapter{Hepatic Vein Thrombosis}

\section*{Definition and Overview}
Hepatic vein thrombosis, also known as Budd--Chiari syndrome, is a rare condition in which blood clots form in the veins that drain blood from the liver. The blockage raises pressure within the liver, causing congestion, liver damage, and ultimately liver failure. It can occur suddenly or develop gradually, and ranges from a silent finding to a life-threatening emergency. The condition is strongly associated with blood clotting disorders and is managed by specialist liver teams using blood thinners, procedures to restore blood flow, and treatment of the underlying cause.

\section*{Epidemiology}
Budd--Chiari syndrome is rare, affecting around one in 100,000 to one in 1,000,000 people per year. It affects all ages but is most common in adults between 20 and 40, with a slight female predominance. Up to half of cases have no identifiable cause, while the remainder are associated with inherited or acquired blood clotting disorders, pregnancy, oral contraceptives, and conditions such as polycythaemia vera. Because it is rare, diagnosis is often delayed, and it should be considered in anyone with unexplained liver dysfunction and clotting risk.

\section*{Aetiology and Pathophysiology}
\begin{itemize}
    \item \textbf{Clotting disorders:} myeloproliferative diseases such as polycythaemia vera and essential thrombocythaemia are the most common causes
    \item \textbf{Inherited thrombophilia:} factor V Leiden, protein C and S deficiency, and antithrombin deficiency
    \item \textbf{Hormonal factors:} oral contraceptives and pregnancy increase clotting risk
    \item \textbf{Other causes:} paroxysmal nocturnal haemoglobinuria, Behçet's disease, cancer, and abdominal trauma
    \item \textbf{Pathophysiology:} obstruction of the hepatic veins raises pressure in the liver's blood sinusoids, causing congestion, swelling, and ischaemic injury; chronic congestion leads to fibrosis and cirrhosis
    \item \textbf{Acute vs chronic:} acute thrombosis presents abruptly; chronic disease develops insidiously with portal hypertension
\end{itemize}

\section*{Clinical Features}
\begin{itemize}
    \item Abdominal pain, especially in the upper right side
    \item Abdominal swelling from fluid accumulation (ascites)
    \item Enlarged, tender liver
    \item Nausea, vomiting, and loss of appetite
    \item Jaundice: yellowing of the skin and eyes
    \item Leg swelling and distended abdominal veins from portal hypertension
    \item Chronic cases present with cirrhosis-like features: varices, bleeding, and hepatic encephalopathy
    \item Fulminant cases cause rapid liver failure with confusion and bleeding
\end{itemize}

\section*{Differential Diagnosis}
\begin{itemize}
    \item Cirrhosis from other causes, including alcohol and viral hepatitis
    \item Congestive heart failure, which causes liver congestion
    \item Veno-occlusive disease (sinusoidal obstruction syndrome) after transplantation or chemotherapy
    \item Gallbladder disease and other causes of right upper quadrant pain
    \item Acute hepatitis from viruses or drugs
    \item Liver tumours obstructing the veins
    \item Imaging of the hepatic veins distinguishes these conditions
\end{itemize}

\section*{When to Seek Medical Care}
Anyone with unexplained abdominal swelling, an enlarged liver, or new liver problems, particularly with a history of blood clots or a clotting disorder, needs urgent specialist assessment.

\textbf{Call 000 (or local emergency services) for:}
\begin{itemize}
    \item Sudden severe abdominal pain with swelling
    \item Vomiting blood or passing black stools (bleeding varices)
    \item Confusion, drowsiness, or personality change (hepatic encephalopathy)
    \item Severe breathlessness or sudden deterioration
    \item Signs of severe bleeding or collapse
\end{itemize}

\section*{Diagnosis and Investigations}
\begin{itemize}
    \item \textbf{Imaging:} Doppler ultrasound is the first-line test, showing absent or reduced flow in the hepatic veins
    \item \textbf{CT and MRI venography:} confirm the diagnosis and map the extent of thrombosis
    \item \textbf{Liver biopsy:} may be used when imaging is inconclusive or chronic disease is suspected
    \item \textbf{Blood tests:} liver function, clotting studies, and tests for underlying clotting disorders and myeloproliferative disease
    \item \textbf{Bone marrow assessment:} for suspected myeloproliferative causes
    \item \textbf{Genetic testing:} for inherited thrombophilia
\end{itemize}

\section*{Management}
\begin{itemize}
    \item \textbf{Anticoagulation:} blood thinners such as warfarin or low-molecular-weight heparin to prevent clot extension and treat the underlying cause
    \item \textbf{Treatment of the underlying disorder:} managing myeloproliferative disease and other causes
    \item \textbf{Interventional radiology:} angioplasty and stenting of narrowed veins, or thrombolysis for acute clots
    \item \textbf{Transjugular intrahepatic portosystemic shunt (TIPS):} creates a new pathway for blood to bypass the blocked liver veins
    \item \textbf{Ascites management:} diuretics and fluid restriction
    \item \textbf{Variceal management:} endoscopic treatment and beta-blockers for bleeding risk
    \item \textbf{Liver transplantation:} for patients with fulminant or end-stage liver failure
\end{itemize}

\section*{Complications}
\begin{itemize}
    \item Progressive liver failure and cirrhosis
    \item Bleeding from oesophageal varices
    \item Ascites and its complications, including spontaneous bacterial peritonitis
    \item Hepatic encephalopathy: confusion and altered consciousness from liver dysfunction
    \item Kidney failure in advanced disease
    \item Recurrent thrombosis despite treatment
    \item Complications of anticoagulation, including bleeding
\end{itemize}

\section*{Prognosis and Outcomes}
The outlook depends on how quickly the condition is recognised and the underlying cause. With anticoagulation and interventional treatment, many patients live for years with good liver function. Acute, fulminant disease carries high mortality without urgent treatment, and transplantation is life-saving for those who progress. Chronic disease is managed like cirrhosis, with outcomes depending on complications. Specialist management and treatment of the clotting disorder are the keys to a good prognosis.

\section*{Prevention and Self-Care}
\begin{itemize}
    \item People with known clotting disorders should discuss thrombosis risk and prevention with their specialist
    \item Consider the risks and benefits of oral contraceptives in women with thrombophilia
    \item Report new abdominal pain, swelling, or jaundice promptly
    \item Attend specialist follow-up and take anticoagulation exactly as prescribed
    \item Avoid alcohol to protect the liver
    \item Keep vaccinations up to date, including hepatitis A and B
    \item Seek genetic counselling if an inherited clotting disorder is found
\end{itemize}

\section*{Sources and Further Reading}
The following reputable sources provide further information for healthcare students and patients seeking reliable, current advice about hepatic vein thrombosis.

\begin{itemize}
    \item Healthdirect. \textit{Liver disease and hepatic vein thrombosis.} https://www.healthdirect.gov.au/
    \item Liver Foundation. \textit{Budd--Chiari syndrome information.} https://liver.org.au/
    \item MSD Manual. \textit{Budd--Chiari syndrome.} https://www.msdmanuals.com/
    \item NHS. \textit{Budd--Chiari syndrome.} https://www.nhs.uk/
    \item American Liver Foundation. \textit{Budd--Chiari syndrome.} https://liverfoundation.org/
    \item MedlinePlus. \textit{Budd-Chiari syndrome.} https://medlineplus.gov/ency/article/000239.htm
\end{itemize}
